\slajdid{09_3-s062}
\begin{frame}[label=09_3-s062]
\gusttekst\setlength{\abovedisplayskip}{3pt}\setlength{\belowdisplayskip}{3pt}\setlength{\abovedisplayshortskip}{2pt}\setlength{\belowdisplayshortskip}{2pt}Slučaj $V_C=+15V,V_I=-12V$\begin{center}\input{slike/tikz/s062_bilateralni_prekidac.tex}\end{center}Sada P tranzistor ima sjajne uslove za provođenje. Sors mu je sa strane 1, pa je napon između sorsa i gejta velik. Zbog toga će radi sa velikom strujom koja je u smeru kako je prikazano na slici. Napon na drejnu će biti pozitivan i on će raditi u omskoj oblasti. Ako je otpornost $R_L$ velika i veća od dinamičke otpornosti tranzistora izlazni napon će biti $V_O\approx+12V$. Bitno je da uočite šta se dešava sa N tranzistorom. Njegov napon između gejta i sorsa je $V_{GS}>V_{Tn}$. Sors mu je na strani 2. Zbog malog napona između drejna i sorsa (obezbeđuje P tranzistor) on je takođe u omskoj oblasti, ali zbog relativno malog napona između gejta i sorsa sa velikom otpornošću.
\end{frame}
